\begingroup

\fontsize{9.4pt}{12.2pt}\selectfont

\setlength{\emergencystretch}{2em}

\sloppy

\begin{thebibliography}{99}

\setlength{\itemsep}{0.25ex}

\bibitem{ref1} Scrosati B, Garche J. Lithium batteries: Status, prospects and future. Journal of Power Sources 2010;195(9):2419–2430. \url{https://doi.org/10.1016/j.jpowsour.2009.11.048}.

\bibitem{ref2} Kim T, Song W, Son D-Y, Ono LK, Qi Y. Lithium-ion batteries: Outlook on present, future, and hybridized technologies. Journal of Materials Chemistry A 2019;7(7):2942–2964. \url{https://doi.org/10.1039/C8TA10513H}.

\bibitem{ref3} Krause FC, Ruiz JP, Jones SC, Brandon EJ, Darcy EC, Iannello CJ, Bugga RV. Performance of commercial Li-ion cells for future NASA missions and aerospace applications. Journal of The Electrochemical Society 2021;168(4):040504. \url{https://doi.org/10.1149/1945-7111/abf05f}.

\bibitem{ref4} Xiong R, Kim J, Shen W, et al. Key technologies for electric vehicles. Green Energy and Intelligent Transportation 2022;1(2):100041. \url{https://doi.org/10.1016/j.geits.2022.100041}.

\bibitem{ref5} Barré A, Deguilhem B, Grolleau S, Gérard M, Suard F, Riu D. A review on lithium-ion battery ageing mechanisms and estimations for automotive applications. Journal of Power Sources 2013;241:680–689. \url{https://doi.org/10.1016/j.jpowsour.2013.05.040}.

\bibitem{ref6} Ma S, Jiang M, Tao P, Song C, Wu J, Wang J, Deng T, Shang W. Temperature effect and thermal impact in lithium-ion batteries: A review. Progress in Natural Science: Materials International 2018;28(6):653–666. \url{https://doi.org/10.1016/j.pnsc.2018.11.002}.

\bibitem{ref7} Wang C-J, Zhu Y-L, Gao F, Bu X-Y, Chen H-S, Quan T, Xu Y-B, Jiao Q-J. Internal short circuit and thermal runaway evolution mechanism of fresh and retired Li-ion batteries with LiFePO4 cathode during overcharge. Applied Energy 2022;328:120224. \url{https://doi.org/10.1016/j.apenergy.2022.120224}.

\bibitem{ref8} Feng L, Jiang L, Liu J, Wang Z, Wei Z, Wang Q. Dynamic overcharge investigations of lithium ion batteries with different state of health. Journal of Power Sources 2021;507:230262. \url{https://doi.org/10.1016/j.jpowsour.2021.230262}.

\bibitem{ref9} Xia Y, Chen G, Zhou Q, Shi X, Shi F. Failure behaviours of 100\% SOC lithium-ion battery modules under different impact loading conditions. Engineering Failure Analysis 2017;82:149–160. \url{https://doi.org/10.1016/j.engfailanal.2017.09.003}.

\bibitem{ref10} Che Y, Hu X, Lin X, Guo J, Teodorescu R. Health prognostics for lithium-ion batteries: Mechanisms, methods, and prospects. Energy \& Environmental Science 2023;16:338–371. \url{https://doi.org/10.1039/D2EE03019E}.

\bibitem{ref11} Ge M-F, Liu Y, Jiang X, Liu J. A review on state of health estimations and remaining useful life prognostics of lithium-ion batteries. Measurement 2021;174:109057. \url{https://doi.org/10.1016/j.measurement.2021.109057}.

\bibitem{ref12} Tian H, Qin P, Li K, Zhao Z. A review of the state of health for lithium-ion batteries: Research status and suggestions. Journal of Cleaner Production 2020;261:120813. \url{https://doi.org/10.1016/j.jclepro.2020.120813}.

\bibitem{ref13} Lu J, Xiong R, Tian J, et al. Deep learning to estimate lithium-ion battery state of health without additional degradation experiments. Nature Communications 2023;14:2760. \url{https://doi.org/10.1038/s41467-023-38458-w}.

\bibitem{ref14} Luo K, Chen X, Zheng H, Shi Z. A review of deep learning approach to predicting the state of health and state of charge of lithium-ion batteries. Journal of Energy Chemistry 2022;74:159–173. \url{https://doi.org/10.1016/j.jechem.2022.06.049}.

\bibitem{ref15} Song K, Hu D, Tong Y, Yue X. Remaining life prediction of lithium-ion batteries based on health management: A review. Journal of Energy Storage 2023;57:106193. \url{https://doi.org/10.1016/j.est.2022.106193}.

\bibitem{ref16} Khaleghi S, Hosen M, Karimi D, et al. Developing an online data-driven approach for prognostics and health management of lithium-ion batteries. Applied Energy 2022;308:118348. \url{https://doi.org/10.1016/j.apenergy.2021.118348}.

\bibitem{ref17} Hu X, Xu L, Lin X, Pecht M. Battery lifetime prognostics. Joule 2020;4:310–346. \url{https://doi.org/10.1016/j.joule.2019.11.018}.

\bibitem{ref18} He W, Williard N, Osterman M, Pecht M. Prognostics of lithium-ion batteries based on Dempster–Shafer theory and the Bayesian Monte Carlo method. Journal of Power Sources 2011;196(23):10314–10321. \url{https://doi.org/10.1016/j.jpowsour.2011.08.040}.

\bibitem{ref19} Wang Y, Tian J, Sun Z, Wang L, Xu R, Li M, Chen Z. A comprehensive review of battery modeling and state estimation approaches for advanced battery management systems. Renewable and Sustainable Energy Reviews 2020;131:110015. \url{https://doi.org/10.1016/j.rser.2020.110015}.

\bibitem{ref20} Deng Z, Yang L, Deng H, Cai Y, Li D. Polynomial approximation pseudo-two-dimensional battery model for online application in embedded battery management system. Energy 2018;142:838–850. \url{https://doi.org/10.1016/j.energy.2017.10.097}.

\bibitem{ref21} Li J, Adewuyi K, Lotfi N, Landers RG, Park J. A single particle model with chemical/mechanical degradation physics for lithium ion battery state of health (SOH) estimation. Applied Energy 2018;212:1178–1190. \url{https://doi.org/10.1016/j.apenergy.2018.01.011}.

\bibitem{ref22} Wang D, Zhang Q, Huang H, Yang B, Dong H, Zhang J. An electrochemical–thermal model of lithium-ion battery and state of health estimation. Journal of Energy Storage 2022;47:103528. \url{https://doi.org/10.1016/j.est.2021.103528}.

\bibitem{ref23} Rechkemmer SK, Zang X, Zhang W, Sawodny O. Empirical Li-ion aging model derived from single particle model. Journal of Energy Storage 2019;21:773–786. \url{https://doi.org/10.1016/j.est.2019.01.005}.

\bibitem{ref24} Chen L, Lü Z, Lin W, Li J, Pan H. A new state-of-health estimation method for lithium-ion batteries through the intrinsic relationship between ohmic internal resistance and capacity. Measurement 2018;116:586–595. \url{https://doi.org/10.1016/j.measurement.2017.11.016}.

\bibitem{ref25} Roman D, Saxena S, Robu V, Flynn D, Pecht M. Machine learning pipeline for battery state of health estimation. arXiv 2021. \url{https://doi.org/10.48550/arXiv.2102.00837}.

\bibitem{ref26} Fei Z, Yang F, Tsui K-L, Li L, Zhang Z. Early prediction of battery lifetime via a machine learning based framework. Energy 2021;225:120205. \url{https://doi.org/10.1016/j.energy.2021.120205}.

\bibitem{ref27} Li Y, Zou C, Berecibar M, Nanini-Maury E, Chan JC-W, van den Bossche P, Van Mierlo J, Omar N. Random forest regression for online capacity estimation of lithium-ion batteries. Applied Energy 2018;232:197–210. \url{https://doi.org/10.1016/j.apenergy.2018.09.182}.

\bibitem{ref28} Richardson RR, Birkl CR, Osborne MA, Howey DA. Gaussian process regression for in situ capacity estimation of lithium-ion batteries. IEEE Transactions on Industrial Informatics 2018;15:127–138. \url{https://doi.org/10.1109/TII.2018.2794997}.

\bibitem{ref29} Qian C, Xu B, Chang L, Sun B, Feng Q, Yang D, Ren Y, Wang Z. Convolutional neural network based capacity estimation using random segments of the charging curves for lithium-ion batteries. Energy 2021;227:120333. \url{https://doi.org/10.1016/j.energy.2021.120333}.

\bibitem{ref30} Lee G, Kwon D, Lee C. A convolutional neural network model for SOH estimation of Li-ion batteries with physical interpretability. Mechanical Systems and Signal Processing 2023;188:110004. \url{https://doi.org/10.1016/j.ymssp.2022.110004}.

\bibitem{ref31} Li X, Zhao M, Zhong S, Li J, Fu S, Yan Z. BMSFormer: An efficient deep learning model for online state-of-health estimation of lithium-ion batteries under high-frequency early SOC data with strong correlated single health indicator. Energy 2024;313:134030. \url{https://doi.org/10.1016/j.energy.2024.134030}.

\bibitem{ref32} Jia C, Tian Y, Shi Y, Jia J, Wen J, Zeng J. State of health prediction of lithium-ion batteries based on bidirectional gated recurrent unit and transformer. Energy 2023;285:129401. \url{https://doi.org/10.1016/j.energy.2023.129401}.

\bibitem{ref33} Zheng Y, Hu J, Chen J, Deng H, Hu W. State of health estimation for lithium battery random charging process based on CNN-GRU method. Energy Reports 2023;9:1–10. \url{https://doi.org/10.1016/j.egyr.2022.12.093}.

\bibitem{ref34} Vaswani A, Shazeer N, Parmar N, Uszkoreit J, Jones L, Gomez AN, Kaiser L, Polosukhin I. Attention Is All You Need. In: Advances in Neural Information Processing Systems, vol. 30; 2017. \url{https://doi.org/10.48550/arXiv.1706.03762}.

\bibitem{ref35} Gu X, See KW, Li P, Shan K, Wang Y, Zhao L, Lim KC, Zhang N. A novel state-of-health estimation for the lithium-ion battery using a convolutional neural network and transformer model. Energy 2023;262:125501. \url{https://doi.org/10.1016/j.energy.2022.125501}.

\bibitem{ref36} Bai T, Wang H. Convolutional transformer-based multiview information perception framework for lithium-ion battery state-of-health estimation. IEEE Transactions on Instrumentation and Measurement 2023;72:1–12. \url{https://doi.org/10.1109/TIM.2023.3300451}.

\bibitem{ref37} Chen L, Xie S, Lopes AM, Bao X. A vision transformer-based deep neural network for state of health estimation of lithium-ion batteries. International Journal of Electrical Power \& Energy Systems 2023;152:109233. \url{https://doi.org/10.1016/j.ijepes.2023.109233}.

\bibitem{ref38} Park SH, Kim SH. Physics-Informed Transformer Using Degradation-Sensitive Indicators for Long-Term State-of-Health Estimation of Lithium-Ion Batteries. Batteries 2026;12(2):48.

\bibitem{ref39} Tay Y, Dehghani M, Bahri D, Metzler D. Efficient Transformers: A survey. ACM Computing Surveys 2022;55(6). \url{https://doi.org/10.1145/3530811}.

\bibitem{ref40} Katharopoulos A, Vyas A, Pappas N, Fleuret F. Transformers are RNNs: Fast autoregressive Transformers with linear attention. In: Proceedings of the 37th International Conference on Machine Learning, PMLR 119; 2020. p. 5156–5165. \url{https://doi.org/10.48550/arXiv.2006.16236}.

\bibitem{ref41} Xu H, Wu L, Xiong S, Li W, Garg A, Gao L. An improved CNN-LSTM model-based state-of-health estimation approach for lithium-ion batteries. Energy 2023;276:127585. \url{https://doi.org/10.1016/j.energy.2023.127585}.

\bibitem{ref42} Tian Y, Wen J, Yang Y, Shi Y, Zeng J. State-of-Health Prediction of Lithium-Ion Batteries Based on CNN-BiLSTM-AM. Batteries 2022;8(10):155. \url{https://doi.org/10.3390/batteries8100155}.

\bibitem{ref43} Liu S, Yu H, Liao C, Li J, Lin W, Liu AX, Dustdar S. Pyraformer: Low-complexity pyramidal attention for long-range time series modeling and forecasting. In: Proceedings of the International Conference on Learning Representations; 2022.

\bibitem{ref44} Zheng L, Wang C, Kong L. Linear complexity randomized self-attention mechanism. In: Proceedings of the 39th International Conference on Machine Learning, PMLR 162; 2022. p. 27011–27041.

\bibitem{ref45} Wang H, Li H, He Z, Chen X, Liu H, Chen C. A linear-complexity deep learning framework for efficient battery state of health estimation in resource-constrained systems. Engineering Applications of Artificial Intelligence 2026;181:115420. \url{https://doi.org/10.1016/j.engappai.2026.115420}.

\bibitem{ref46} Han D, Ye T, Han Y, Xia Z, Pan S, Wan P, Song S, Huang G. Agent Attention: On the integration of Softmax and linear attention. In: Computer Vision—ECCV 2024, Lecture Notes in Computer Science, vol. 15108, Springer, Cham; 2025. p. 124–140. \url{https://doi.org/10.1007/978-3-031-72973-7\_8}.

\bibitem{ref47} Lin C, Xu J, Hou J, et al. A fast data-driven battery capacity estimation method under non-constant current charging and variable temperature. Energy Storage Materials 2023;63:102967. \url{https://doi.org/10.1016/j.ensm.2023.102967}.

\bibitem{ref48} Dai H, Wang J, Huang Y, Lai Y, Zhu L. Lightweight state-of-health estimation of lithium-ion batteries based on statistical feature optimization. Renewable Energy 2024;222:119907. \url{https://doi.org/10.1016/j.renene.2023.119907}.

\bibitem{ref49} Wu J, Liu Z, Zhang Y, Lei D, Zhang B, Cao W. Data-driven state of health estimation for lithium-ion battery based on voltage variation curves. Journal of Energy Storage 2023;73:109191. \url{https://doi.org/10.1016/j.est.2023.109191}.

\bibitem{ref50} Wen J, Chen X, Li X, Li Y. SOH prediction of lithium battery based on IC curve feature and BP neural network. Energy 2022;261:125234. \url{https://doi.org/10.1016/j.energy.2022.125234}.

\bibitem{ref51} Tian Y, Dong Q, Tian J, Li X, Li G, Mehran K. Capacity estimation of lithium-ion batteries based on optimized charging voltage section and virtual sample generation. Applied Energy 2023;332:120516. \url{https://doi.org/10.1016/j.apenergy.2022.120516}.

\bibitem{ref52} Li X, Zhao M, Zhong S, Li J, Cui Z, Fu S, Yan Z. Deep transfer learning enabled online state-of-health estimation of lithium-ion batteries under small samples across different cathode materials, ambient temperature and charge-discharge protocols. Journal of Power Sources 2025;650:237503. \url{https://doi.org/10.1016/j.jpowsour.2025.237503}.

\bibitem{ref53} Tang Q, He T, Huang J, Zhu W. A novel state of health estimation method for lithium-ion battery based on multi-feature extraction and fusion model with bidirectional feature extraction. Journal of Energy Storage 2026;141:119415.

\bibitem{ref54} Gong Y, Lu X, Su Z, Huang Z, Quan Y, Zhang Z, Ouyang T. Correlation-aware fusion model for robust SOH estimation in lithium-ion batteries. Journal of Energy Storage 2026;149:120018.

\bibitem{ref55} Lin M, Ke L, Wang W, Meng J, Guan Y, Wu J. Health prognosis via feature optimization and convolutional neural network for lithium-ion batteries. Engineering Applications of Artificial Intelligence 2024;133:108666. \url{https://doi.org/10.1016/j.engappai.2024.108666}.

\bibitem{ref56} Ambroise C, McLachlan GJ. Selection bias in gene extraction on the basis of microarray gene-expression data. Proceedings of the National Academy of Sciences of the United States of America 2002;99(10):6562–6566. \url{https://doi.org/10.1073/pnas.102102699}.

\bibitem{ref57} Kaufman S, Rosset S, Perlich C, Stitelman O. Leakage in data mining: Formulation, detection, and avoidance. ACM Transactions on Knowledge Discovery from Data 2012;6(4):15. \url{https://doi.org/10.1145/2382577.2382579}.

\bibitem{ref58} Huang K, Yao K, Guo Y, Lv Z. State of health estimation of lithium-ion batteries based on fine-tuning or rebuilding transfer learning strategies combined with new features mining. Energy 2023;282:128739. \url{https://doi.org/10.1016/j.energy.2023.128739}.

\bibitem{ref59} Wang R, Song C, Gao M, Zhao J, Xu Z. Model-data fusion domain adaptation for battery state of health estimation with fewer data and simplified feature extractor. Journal of Energy Storage 2023;60:106686. \url{https://doi.org/10.1016/j.est.2023.106686}.

\bibitem{ref60} Birkl CR, Roberts MR, McTurk E, Bruce PG, Howey DA. Degradation diagnostics for lithium-ion cells. Journal of Power Sources 2017;341:373–386. \url{https://doi.org/10.1016/j.jpowsour.2016.12.011}.

\bibitem{ref61} dos Reis G, Strange C, Yadav M, Li S. Lithium-ion battery data and where to find it. Energy and AI 2021;5:100081. \url{https://doi.org/10.1016/j.egyai.2021.100081}.

\bibitem{ref62} Severson KA, Attia PM, Jin N, et al. Data-driven prediction of battery cycle life before capacity degradation. Nature Energy 2019;4(5):383–391. \url{https://doi.org/10.1038/s41560-019-0356-8}.

\bibitem{ref63} Lin C, Xu J, Shi M, Mei X. Constant current charging time based fast state-of-health estimation for lithium-ion batteries. Energy 2022;247:123556. \url{https://doi.org/10.1016/j.energy.2022.123556}.

\bibitem{ref64} Li D, Xing Y, Li X. Joint state of health and remaining useful life assessment of lithium-ion batteries using a novel attention mechanism-based deep learning model and multi-source optimized health indicators. Journal of Energy Storage 2026;161:121823.

\bibitem{ref65} Wu B, Yufit V, Merla Y, Martinez-Botas RF, Brandon NP, Offer GJ. Differential thermal voltammetry for tracking of degradation in lithium-ion batteries. Journal of Power Sources 2015;273:495–501. \url{https://doi.org/10.1016/j.jpowsour.2014.09.127}.

\bibitem{ref66} Lin Z, Li D, Zou Y. Energy efficiency of lithium-ion batteries: Influential factors and long-term degradation. Journal of Energy Storage 2023;74:109386. \url{https://doi.org/10.1016/j.est.2023.109386}.

\bibitem{ref67} Bai J, Huang J, Luo K, Yang F, Xian Y. A feature reuse based multi-model fusion method for state of health estimation of lithium-ion batteries. Journal of Energy Storage 2023;70:107965. \url{https://doi.org/10.1016/j.est.2023.107965}.

\bibitem{ref68} Zhang Y, Zhang D, Wu T. Method for estimating the state of health of lithium-ion batteries based on differential thermal voltammetry and sparrow search algorithm-Elman neural network. Energy Engineering 2025;122(1):203–220. \url{https://doi.org/10.32604/ee.2024.056244}.

\bibitem{ref69} Zou B, Wang H, Zhang T, et al. A deep learning approach for state-of-health estimation of lithium-ion batteries based on differential thermal voltammetry and attention mechanism. Frontiers in Energy Research 2023;11:1178151. \url{https://doi.org/10.3389/fenrg.2023.1178151}.

\bibitem{ref70} Feng H, Zhang L. A heterogeneous learner fusion method with supplementary feature for lithium-ion batteries state of health estimation. Journal of Energy Storage 2024;92:111896. \url{https://doi.org/10.1016/j.est.2024.111896}.

\bibitem{ref71} Chollet F. Xception: Deep learning with depthwise separable convolutions. In: Proceedings of the IEEE Conference on Computer Vision and Pattern Recognition; 2017. p. 1800–1807. \url{https://doi.org/10.1109/CVPR.2017.195}.

\bibitem{ref72} Sandler M, Howard A, Zhu M, Zhmoginov A, Chen L-C. MobileNetV2: Inverted residuals and linear bottlenecks. In: Proceedings of the IEEE/CVF Conference on Computer Vision and Pattern Recognition; 2018. p. 4510–4520. \url{https://doi.org/10.1109/CVPR.2018.00474}.

\bibitem{ref73} Zhang X, Zhou X, Lin M, Sun J. ShuffleNet: An extremely efficient convolutional neural network for mobile devices. arXiv 2017. \url{https://doi.org/10.48550/arXiv.1707.01083}.

\bibitem{ref74} Ma N, Zhang X, Zheng H-T, Sun J. ShuffleNet V2: Practical guidelines for efficient CNN architecture design. In: Computer Vision—ECCV 2018, Lecture Notes in Computer Science, vol. 11218, Springer, Cham; 2018. p. 116–131. \url{https://doi.org/10.1007/978-3-030-01264-9\_8}.

\bibitem{ref75} Srivastava N, Hinton G, Krizhevsky A, Sutskever I, Salakhutdinov R. Dropout: A simple way to prevent neural networks from overfitting. Journal of Machine Learning Research 2014;15(56):1929–1958.

\bibitem{ref76} Chicco D, Warrens MJ, Jurman G. The coefficient of determination R-squared is more informative than SMAPE, MAE, MAPE, MSE and RMSE in regression analysis evaluation. PeerJ Computer Science 2021;7:e623. \url{https://doi.org/10.7717/peerj-cs.623}.

\bibitem{ref77} Li P, Zhang Z, Grosu R, Deng Z, Hou J, Rong Y, Wu R. An end-to-end neural network framework for state-of-health estimation and remaining useful life prediction of electric vehicle lithium batteries. Renewable and Sustainable Energy Reviews 2022;156:111843. \url{https://doi.org/10.1016/j.rser.2021.111843}.

\end{thebibliography}

\endgroup